\begin{abstract}
Pre-screening patients for clinical trials means checking a free-text note against every eligibility criterion
and citing the facts that justify each decision. We present \sys, a locally executable system for
criterion-level eligibility prediction with evidence tracing, built on Qwen2.5-1.5B-Instruct and the TrialGPT
criterion-annotation benchmark. We fine-tune the model with LoRA on a 6\,GB consumer GPU, choosing bf16 with
gradient checkpointing after measuring the alternatives. On held-out test patients, the same model goes from
0.160 to 0.733 accuracy, from 0.155 to 0.528 macro-F1 over the labels present, and from 0.169 to 0.561
evidence micro-F1, with no unparseable outputs. We also implement and benchmark three inference mechanisms by
hand:
\begin{itemize}
\item a KV-cached decoder, $3.6\times$ faster than recomputation at 128 context tokens and up to $20\times$ at
1{,}024;
\item a tiled online-softmax attention kernel that matches a float64 reference to $2.5\times10^{-6}$ in fp32 and uses
46$\times$ less memory than unfused attention at 4{,}096 tokens;
\item a greedy speculative decoder whose draft acceptance reaches 91\% (96\% with a fine-tuned draft), but
which is still slower than plain cached decoding on launch-bound consumer hardware.
\end{itemize}
Where bf16 outputs diverge from the reference, we trace the cause to logit ties within one bf16 rounding
step and confirm exact agreement in fp32. A browser application screens unseen notes, highlights evidence
sentences, explains each decision only from the cited evidence, and exposes the inference mechanisms live.
\end{abstract}

\begin{IEEEkeywords}
clinical-trial matching, LoRA, KV cache, speculative decoding, online softmax, evidence attribution
\end{IEEEkeywords}
